% Options for packages loaded elsewhere
\PassOptionsToPackage{unicode}{hyperref}
\PassOptionsToPackage{hyphens}{url}
\PassOptionsToPackage{dvipsnames,svgnames,x11names}{xcolor}
\documentclass[
  10pt,
  fontset=fandol]{ctexart}
\usepackage{xcolor}
\usepackage[margin=16mm]{geometry}
\usepackage{amsmath,amssymb}
\setcounter{secnumdepth}{-\maxdimen} % remove section numbering
\usepackage{iftex}
\ifPDFTeX
  \usepackage[T1]{fontenc}
  \usepackage[utf8]{inputenc}
  \usepackage{textcomp} % provide euro and other symbols
\else % if luatex or xetex
  \usepackage{unicode-math} % this also loads fontspec
  \defaultfontfeatures{Scale=MatchLowercase}
  \defaultfontfeatures[\rmfamily]{Ligatures=TeX,Scale=1}
\fi
\usepackage{lmodern}
\ifPDFTeX\else
  % xetex/luatex font selection
\fi
% Use upquote if available, for straight quotes in verbatim environments
\IfFileExists{upquote.sty}{\usepackage{upquote}}{}
\IfFileExists{microtype.sty}{% use microtype if available
  \usepackage[]{microtype}
  \UseMicrotypeSet[protrusion]{basicmath} % disable protrusion for tt fonts
}{}
\makeatletter
\@ifundefined{KOMAClassName}{% if non-KOMA class
  \IfFileExists{parskip.sty}{%
    \usepackage{parskip}
  }{% else
    \setlength{\parindent}{0pt}
    \setlength{\parskip}{6pt plus 2pt minus 1pt}}
}{% if KOMA class
  \KOMAoptions{parskip=half}}
\makeatother
\setlength{\emergencystretch}{3em} % prevent overfull lines
\providecommand{\tightlist}{%
  \setlength{\itemsep}{0pt}\setlength{\parskip}{0pt}}
\usepackage{amsmath,amssymb,mathtools,needspace}
\xeCJKDeclareCharClass{CJK}{"2032,"0400->"04FF,"2160->"216B,"2460->"2473,"25A0,"25C7,"25CB,"25CF}
\IfFontExistsTF{Arial Unicode MS}{\xeCJKsetup{AutoFallBack=true}\setCJKfallbackfamilyfont{\CJKrmdefault}{Arial Unicode MS}}{}
\setcounter{secnumdepth}{0}
\setlength{\emergencystretch}{2em}
\allowdisplaybreaks
\usepackage{bookmark}
\IfFileExists{xurl.sty}{\usepackage{xurl}}{} % add URL line breaks if available
\urlstyle{same}
\hypersetup{
  pdftitle={三次样条插值的收敛性},
  colorlinks=true,
  linkcolor={Maroon},
  filecolor={Maroon},
  citecolor={Blue},
  urlcolor={Blue},
  pdfcreator={LaTeX via pandoc}}

\title{三次样条插值的收敛性}
\author{}
\date{}

\begin{document}
\maketitle

\subsubsection{2.8.5 三次样条插值的收敛性}\label{ux4e09ux6b21ux6837ux6761ux63d2ux503cux7684ux6536ux655bux6027}

为了证明三次样条插值的收敛性，需要用到向量与矩阵范数及有关结论，可参看 7.5 节。

设 \(\boldsymbol A=(a_{ij})_n\) 为 \(n\times n\) 矩阵，\(\boldsymbol x=(x_1,x_2,\cdots,x_n)^{\mathrm T}\) 为 \(n\) 维向量，定义 \(\boldsymbol x\) 及 \(\boldsymbol A\) 的\textbf{范数}为

\[\begin{gathered}
\|\boldsymbol x\|_\infty=\max_{1\leqslant i\leqslant n}|x_j|,\\
\|\boldsymbol A\|_\infty=\sup_{\boldsymbol x\ne0}\frac{\|\boldsymbol A\boldsymbol x\|_\infty}{\|\boldsymbol x\|_\infty}.
\end{gathered}\tag{2.8.21}\]

同样，对于函数 \(f(x)\in C[a,b]\)，也定义 \(f\) 的范数为

\[\|f\|_\infty=\sup_{a\leqslant x\leqslant b}|f(x)|=\max_{a\leqslant x\leqslant b}|f(x)|.\]

\textbf{引理}　若 \(\boldsymbol A=(a_{ij})_n\) 具有强对角占优，即

\[\sum_{j\ne i}|a_{ij}|<|a_{ii}|\quad(i=1,2,\cdots,n),\tag{2.8.22}\]

则 \(\boldsymbol A^{-1}\) 存在，且

\[\|\boldsymbol A^{-1}\|_\infty\leqslant\left\{\min\left(|a_{ij}|-\sum_{j\ne i}|a_{ij}|\right)\right\}^{-1}.\tag{2.8.23}\]

该引理的证明可参考定理 8.6 和第 8 章习题 20。

现在以自然边界条件（2.8.4）′的三次样条插值函数 \(S(x)\) 为例，讨论其收敛性，

\begin{quote}
续页：收敛性定理见下一页。
\end{quote}

\begin{quote}
\textbf{校注（agent 补充；原书疑误）}：式（2.8.21）原图在 \(\max\) 下用 \(i\)，而被取最大值的分量印为 \(x_j\)；已原样保留，标准表达需将哑指标统一。式（2.8.23）括号内第一项原图印为 \(|a_{ij}|\)；结合（2.8.22）的对角占优量，应为 \(|a_{ii}|\)，并对行指标 \(i\) 取最小值。参见\href{../assets/ch02b/pdf-054-norm-detail.png}{范数公式原图裁切}。
\end{quote}

\begin{quote}
\textbf{校注（agent 补充；原书表值与条件不一致）}：表 2.8 的 \(S_{10}\) 数值在此按原图保留。按上一页明定的节点、\(f(x)=1/(1+x^2)\) 及 \(S'_{10}(\pm5)=f'(\pm5)\)，精确有理数解三转角方程给出 \(S_{10}(-4.8)\approx0.04162183\)、\(S_{10}(-4.5)\approx0.04716801\)，与原表的 \(0.03758\)、\(0.04248\) 不同；这不是 OCR 差异。复核结果见\href{../../staging/ch02b/table-2-8-check.txt}{独立数值检查}。本校注不以复算值替换原表。
\end{quote}

\href{../../source-images/pdf-055.jpeg}{查看原页}

\begin{quote}
本页接上页 2.8.5 正文。
\end{quote}

此时方程为式（2.8.14），写成

\[\boldsymbol A\boldsymbol m=\boldsymbol g,\tag{2.8.24}\]

其中

\[\boldsymbol A=\begin{pmatrix}
2&1&0&\cdots&0\\
\lambda_1&2&\mu_1&\ddots&\vdots\\
0&\lambda_2&\ddots&\ddots&0\\
\vdots&\ddots&\ddots&2&\mu_{n-1}\\
0&\cdots&0&\lambda_n&2
\end{pmatrix},\quad
\boldsymbol m=\begin{pmatrix}m_0\\m_2\\\vdots\\m_{n-1}\\m_n\end{pmatrix},\quad
\boldsymbol g=\begin{pmatrix}g_0\\g_1\\\vdots\\g_{n-1}\\g_n\end{pmatrix}.\]

由于 \(\mu_i+\lambda_i=1\)（\(i=0,1,\cdots,n\)），且 \(a_{ii}=2\)，故由引理得

\[\|\boldsymbol A^{-1}\|\leqslant1.\tag{2.8.25}\]

\textbf{定理 2.3}　若 \(f(x)\in C[a,b]\)，\(S(x)\) 是以 \(a=x_0<x_1<\cdots<x_n=b\) 为节点，满足条件式（2.8.1）及式（2.8.4）′的三次样条插值函数，令

\[h_j=x_{j+1}-x_j,\quad h=\max_{0\leqslant j\leqslant n-1}h_j,\quad\delta=\min_{0\leqslant j\leqslant n-1}h_j,\]

设 \(h/\delta\leqslant c<\infty\)，则 \(S(x)\) 在 \([a,b]\) 上一致收敛到 \(f(x)\)。

\textbf{证明}　由于 \(S(x)\) 可用式（2.8.6）表示，当 \(x\in[a,b]\) 时，利用式（2.7.9）可得到

\[\|f(x)-S(x)\|_\infty\leqslant\omega(h)+(8/27)h\|\boldsymbol m\|_\infty,\tag{2.8.26}\]

由方程（2.8.24），并利用式（2.8.25）得

\[\|\boldsymbol m\|_\infty\leqslant\|\boldsymbol A^{-1}\|_\infty\|\boldsymbol g\|_\infty\leqslant\|\boldsymbol g\|_\infty^m,\tag{2.8.27}\]

根据式（2.8.11）及式（2.8.13）′可估得

\[\|\boldsymbol g\|_\infty\leqslant3\max_{0\leqslant j\leqslant n-1}|f[x_j,x_{j+1}]|\leqslant\frac3\delta\omega(h)\leqslant\frac6\delta\|f\|_\infty,\tag{2.8.28}\]

将式（2.8.9）及式（2.8.28）代入式（2.8.26），得

\[\|f(x)-S(x)\|_\infty\leqslant\left(1+\frac89\frac h\delta\right)\omega(h).\]

由于 \(f(x)\in C[a,b]\)，当 \(h\to0\) 时，\(\|f(x)-S(x)\|_\infty\to0\)，故 \(S(x)\) 在 \([a,b]\) 上一致收敛到 \(f(x)\)。证毕。

如果 \(f(x)\in C^1[a,b]\)，则 \(S'(x)\) 在 \([a,b]\) 上也一致收敛到 \(f'(x)\)。证明从略。

其他边界条件的三次样条插值函数收敛性证明与此类似，读者可自行证明。

\begin{center}\rule{0.5\linewidth}{0.5pt}\end{center}

\subsection{相关原页校注（agent 补充）}\label{ux76f8ux5173ux539fux9875ux6821ux6ce8agent-ux8865ux5145}

以下校注来自本节覆盖的原页，可能也涉及同页相邻小节；原文照录与校正说明分开保留。

\subsubsection{原书 PDF 页 55 的校注}\label{ux539fux4e66-pdf-ux9875-55-ux7684ux6821ux6ce8}

\href{../pages/pdf-055.md}{完整原页转录} · \href{../../source-images/pdf-055.jpeg}{原页图像}

校注记录：ch02b-E009, ch02b-E010, ch02b-E011。

\begin{quote}
\textbf{校注（agent 补充；原书疑误）}：本页（2.8.24）下面矩阵的末行次对角元印为 \(\lambda_n\)，而自然边界的（2.8.14）对应位置为 \(1\)；向量 \(\boldsymbol m\) 的第二项印为 \(m_2\)，按方程顺序应为 \(m_1\)。上述均保留原图，参见\href{../assets/ch02b/pdf-055-matrix-detail.png}{矩阵与向量原图裁切}。本页没有为自然边界另定义 \(\lambda_n\)，不能沿用前页周期边界的同名系数。
\end{quote}

\begin{quote}
\textbf{校注（agent 补充；原书疑误）}：式（2.8.27）最后一个范数原图上方确有多余的 \(m\)，已按原式转录为 \(\|\boldsymbol g\|_\infty^m\)；由（2.8.25）应得 \(\|\boldsymbol m\|_\infty\leqslant\|\boldsymbol g\|_\infty\)，无这个幂次。参见\href{../assets/ch02b/pdf-055-eq-2-8-27-detail.png}{原式放大裁图}。其后``将式（2.8.9）及式（2.8.28）代入''中的（2.8.9）应指此处的范数估计（2.8.27）；原文编号已保留。
\end{quote}

\subsection{本节覆盖原页的校注记录}\label{ux672cux8282ux8986ux76d6ux539fux9875ux7684ux6821ux6ce8ux8bb0ux5f55}

下列记录按原页关联，含同页相邻小节的校注；计算时同时读取上文校注和完整原页。

\begin{itemize}
\tightlist
\item
  \texttt{ch02b-E006}：原书 PDF 页 54
\item
  \texttt{ch02b-E007}：原书 PDF 页 54
\item
  \texttt{ch02b-E008}：原书 PDF 页 54
\item
  \texttt{ch02b-E009}：原书 PDF 页 55
\item
  \texttt{ch02b-E010}：原书 PDF 页 55
\item
  \texttt{ch02b-E011}：原书 PDF 页 55
\end{itemize}

\end{document}
